\section{Selective seed expansion and exact basis completion}
\label{app:basis-completion}
The seed-expansion study asks whether a smaller first calculation and an exact conversion of its remaining integrals can be cheaper than one large reduction. It uses the seven families in Appendix~\ref{app:workflow} and the target definitions in Appendix~\ref{app:type-specialists}. Development
uses profile $h=71$, training uses $h=72,73$, validation uses $h=74$, and
the reserved test uses $h=75,76$. These are separate rational parameter
profiles. The graph structures and target index patterns occur in all
splits, and $d$ remains symbolic.

\subsection{Changing individual seed bounds}
The bounds $B_c=(r_c,s_c,q_c)$ retain the definitions in
Appendix~\ref{app:workflow}: $r_c$ bounds the sum of positive indices,
$s_c$ bounds the numerator rank, and $q_c$ bounds the number of extra
denominator powers. Each component starts from the largest value among
the requested targets. We first enforce the sector lower bound in
Eq.~\eqref{eq:type-valid-bound}. The six new direct reductions use
\begin{equation}
 B_c+\boldsymbol\delta,\qquad
 \boldsymbol\delta\in\{(1,0,0),(0,1,0),(0,0,1),
                         (1,0,1),(1,1,0),(0,1,1)\}.
 \label{eq:basis-deltas}
\end{equation}
For a seed with $t$ positive indices, summing $a_i=1+(a_i-1)$ over those
indices gives $r=t+q$. The positive-index and extra-power limits can
therefore overlap, and different requested triples can give the same
seed region after all limits are applied. The labels therefore identify requested limits rather than necessarily distinct equation sets. Their purpose is to allow selective expansion, such as requesting more numerator identities without also increasing both denominator limits. Sufficiency remains a property to check for the resulting system. Each direct
reduction keeps the prescribed-basis check described in Ref.~\cite{Kira3}.
If it fails, the calculation calls the preceding fixed graph-specific
method and includes both attempts in its time.

\subsection{Completing a temporary basis}
We also test two ways to complete a smaller reduction. The first stage
uses $B_c$ and supplies the preferred masters, while allowing Kira to
return additional integrals. Kira supports this behaviour when its early
master-basis check is disabled~\cite{Kira3}. Let $M_m$ denote prescribed
masters and let $U_u$ denote additional integrals that remain in the
first-stage output. For target $I_t$, this stage returns exact relations
\begin{equation}
 I_t=\sum_m A_{tm}M_m+\sum_u C_{tu}U_u.
 \label{eq:basis-intermediate}
\end{equation}
Here $A_{tm}$ is the coefficient of prescribed master $M_m$, and
$C_{tu}$ is the coefficient of unresolved integral $U_u$. These
coefficients are rational functions of the kinematics and $d$.
The integrals $U_u$ are unresolved in the smaller first-stage system and need not be independent masters of the full family. Their index vectors define the targets of the second reduction, allowing that calculation to focus on the integrals that actually remain. That request uses its own componentwise target bounds
plus $(1,1,1)$ and again supplies the prescribed basis. One candidate uses
the maximal sectors of the $U_u$ and their subsectors; the other keeps the
original top-level sectors. Both retain the positive-index lower bound.

If the second reduction succeeds, it gives
\begin{equation}
 U_u=\sum_m H_{um}M_m.
 \label{eq:basis-bridge}
\end{equation}
The coefficient $H_{um}$ multiplies master $M_m$ in the reduction of $U_u$.
Substitution of Eq.~\eqref{eq:basis-bridge} into
Eq.~\eqref{eq:basis-intermediate} yields
\begin{align}
 I_t
 &=\sum_m A_{tm}M_m+\sum_u C_{tu}\sum_mH_{um}M_m\\
 &=\sum_m\left(A_{tm}+\sum_uC_{tu}H_{um}\right)M_m.
 \label{eq:basis-compose}
\end{align}
The final coefficient matrix is therefore $A+CH$. Computing its entries exactly in the rational-function field, cancelling common factors, and checking the final integral labels gives a reduction in the prescribed basis. When either reduction fails, the preceding method supplies a full
fallback. If no additional integral occurs, the second reduction is
unnecessary. All basis conversion and fallback costs enter the timer.
Independent reference coefficients check the final result only; the
policy and the two reduction stages do not receive them.

\subsection{Training, model selection, and comparison}
The development set compares the preceding method, six direct expansions,
and two basis-completion choices. We retain the preceding method and each
case's fastest candidate when it lowers that case's complete time by at
least ten percent. This rule selects candidates before training and does
not inspect reserved data.

We then fit the ten model types and nine target features defined in
Appendix~\ref{app:type-specialists}, separately for each graph. The new
response is $\log(t_{cj}/t_{c\mathrm{prev}})$, so it directly measures
improvement over the preceding method. Here $t_{cj}$ is the median complete
time of candidate $j$ on batch $c$, and $t_{c\mathrm{prev}}$ is the
corresponding time of the preceding method. Two repetitions supply each
training label, and three repetitions supply each validation cost. A
fixed strategy minimizes geometric mean training time for each graph.
Validation first selects the faster control between that fixed strategy
and the preceding method. It then tests the fitted predictors at margins
$0,0.05,0.15$ and enables one only when it lowers the control's validation
time by at least five percent. The fitting formulas and their derivations
remain those in Appendices~\ref{app:workflow} and
\ref{app:type-specialists}.

The fixed control may call the learned method retained from the preceding study. Its purpose is to isolate the benefit of the new choice among candidates while preserving the earlier improvements in every method.

The chosen graph policy and candidate set are fixed before the reserved
tests. Inputs with another graph structure or target features outside the
training range use the preceding method. That method retains its own
native fallback. Each reserved batch compares native choices, the
preceding method, the new fixed strategy, and the selected new method on
the same computer. Five repetitions use a rotating order. All methods
use the same executable and exact final coefficient checks. The time
includes preparation, prediction, Kira, any basis conversion, checking,
cleanup, and every retry. Offline training, independent reference
generation, and writing benchmark records remain outside these times.
The selection program and fitted parameters are loaded
before these trials; each attempt starts a fresh Kira process. Separate
application tests start a new caller and supply only the targets, family,
and prescribed basis. They check the coefficients after the reduction.

We first take the median time for each batch and method, then combine
ratios within each graph and give all seven graphs equal weight. The
geometric-mean derivation in Appendix~\ref{app:workflow} also applies here.
Descriptive aggregate intervals resample the seven graph groups. The
per-graph bars show the range over the two reserved parameter profiles.
These ranges describe the tested profiles and do not establish a gain
for all integrals with the same graph.
